\documentclass[a4paper, 11pt]{article}

% ==========================================
% PAQUETES PRINCIPALES (COMPATIBLE OVERLEAF)
% ==========================================
\usepackage[utf8]{inputenc}     % Codificación de caracteres
\usepackage[T1]{fontenc}        % Codificación de fuentes (importante para español)
%\usepackage[spanish, es-tabla]{babel} % Idioma español

% --- Matemáticas y Símbolos ---
\usepackage{amsmath}  % Ecuaciones avanzadas
\usepackage{amssymb}  % Símbolos matemáticos
\usepackage{amsfonts} % Fuentes matemáticas
\usepackage{mathtools} % Mejoras visuales para matemáticas

% --- Estilo y Formato ---
\usepackage[a4paper, top=2.5cm, bottom=2.5cm, left=2.5cm, right=2.5cm]{geometry} % Márgenes
\usepackage{mathptmx} % Usa fuente tipo Times (muy académica)
\usepackage{setspace} % Para interlineado
\onehalfspacing       % Interlineado 1.5 (más legible)

% --- Utilidades ---
\usepackage{listings}
\usepackage{xcolor}
\definecolor{codegreen}{rgb}{0,0.6,0}
\definecolor{codegray}{rgb}{0.5,0.5,0.5}
\definecolor{codepurple}{rgb}{0.58,0,0.82}
\definecolor{backcolour}{rgb}{0.95,0.95,0.92}
\usepackage{graphicx} % Para imágenes
\usepackage{hyperref} % Para enlaces y referencias interactivas
\usepackage{authblk}
\usepackage{enumitem} % Personalización de listas
\newcommand\C[1]\null

% Configuración de enlaces (colores discretos)
\hypersetup{
    colorlinks=true,
    linkcolor=blue,
    filecolor=magenta,      
    urlcolor=cyan,
    citecolor=black
}

%BIBLIOGRAFÍA - PUEDES ELEGIR ENTRE ESTILO IEEE O APA. POR DEFECTO ESTÁ CONFIGURADO IEEE. SI DESEAS USAR APA, COMENTA LAS LÍNEA DE IEEE Y DESCOMENTA LAS DE APA. Si haces cambios en la configuración de la bibliografía y no obtienes los resultados esperados, es recomendable limpiar los archivos auxiliares y volver a compilar en este orden: COMPILAR-BIBLIOGRAFIA-COMPILAR
% Tienes más información sobre cómo generar bibliografía en http://tex.stackexchange.com/questions/154751/biblatex-with-biber-configuring-my-editor-to-avoid-undefined-citations , https://es.sharelatex.com/learn/Bibliography_management_in_LaTeX y en http://www.ctan.org/tex-archive/macros/latex/exptl/biblatex-contrib
% También te recomendamos consultar la guía temática de la Biblioteca sobre citas bibliográficas: http://uc3m.libguides.com/guias_tematicas/citas_bibliograficas/inicio

% CONFIGURACIÓN PARA LA BIBLIOGRAFÍA IEEE
\usepackage[backend=biber, style=ieee, isbn=false,sortcites, maxbibnames=5, minbibnames=1]{biblatex} % Configuración para el estilo de citas de IEEE, recomendado para el área de ingeniería. "maxbibnames" indica que a partir de 5 autores trunque la lista el primero (minbibnames) y añada "et al." tal y como se utiliza en el estilo IEEE.

%CONFIGURACIÓN PARA LA BIBLIOGRAFÍA APA
%\usepackage[style=apa, backend=biber, natbib=true, hyperref=true, uniquelist=false, sortcites]{biblatex}
%\DeclareLanguageMapping{spanish}{spanish-apa}

\addbibresource{bibliografia/bibliografia.bib} % llama al archivo bibliografia.bib que utilizamos de ejemplo




\title{Path Planning and Obstacle Avoidance Real-time Collision Evasion}
\author[1,2]{Pablo De Porcellinis}
\author[1]{Xabier Olaz}
\author[1,3]{Jesús Villadangos}
\affil[1]{Mathematical Engineering and Computer Science Department, Public University of Navarre, Spain}
\affil[2]{FuVeX Civil SL, Spain}
\affil[3]{Institute of Smart Cities, Public University of Navarre,, Spain}

\date{}

\begin{document}

\maketitle

\begin{abstract}

\textbf{Keywords:..}

\end{abstract}

\section{Introduction}

Technological improvement has been present in the modern world. Over the last few decades, the contribution of technologies not only has an impact in the social and political changes but have pervade in every dimension of daily life and industrialized societies \cite{johnson2021technology}.
One of the fields that has been impacted by recent technological development is the military field, as the ongoing global conflicts represent a source of money and a perfect window for technological development. There is a recent example of this approach in the conflict between Ukraine and Russia, which has shown how technology can play an important role in modern warfare, especially accelerating the implementation of fully autonomous solutions, such as Unmanned Aerial Vehicles \cite{samakashvili2024technology}.

Thanks to this development of UAS technology from the military side, it opens a new window also from a civilian perspective, as a transformative technology in many different fields. This impact is highlighted in fields such as agriculture, which offers a significant environmental and economic impact \cite{imran2025environmental}, critical infrastructures, which offer new maintenance inspection mechanisms or surveillance and protection of these national security assets \cite{marut2024surveillance}.

However, this emerging industry around UAS technology is currently at a technological and regulatory juncture. Although the electromechanical and aerodynamic capabilities of modern drones have reached a level of maturity sufficient to reliably perform complex inspections of critical infrastructure—most notably electrical transmission lines—extensive experience in operational power-line inspection reveals that the autonomous control systems governing these platforms remain largely rooted in academic paradigms that fail to capture the constraints and complexities of operational reality. Most state-of-the-art (SOTA) algorithms address geometric optimization problems in static or highly simplified environments, systematically overlooking the dynamic regulatory constraints imposed by frameworks such as EASA SORA 2.5 \cite{EASA_AMC_GM_SORA25_2025}.

This research report addresses the existing gap between engineering capability and regulatory requirements. Its central objective is to conduct a comprehensive review of the state of the art in global and local path planning to assess the feasibility, distinctiveness, and strategic advantage of the proposed Reactive Local Model. The fundamental thesis of this document is that current algorithms—ranging from Artificial Potential Fields (APF) to Rapidly-exploring Random Trees (RRT)—fail by treating safety as a secondary or ‘soft’ optimization variable, rather than as a hard, inviolable constraint derived from aeronautical law and, more precisely, from UAS law.

Taking into account these complexities, the selection of the path-planning method is critical to the operational success of the UAS. Through an in-depth analysis of the academic literature, the current regulatory framework, and available commercial solutions, this report validates the proposal to implement a cost function based on Control Barrier Functions (CBFs) with infinite penalization. This approach constitutes a critical market differentiation: a shift from conventional “obstacle avoidance” (not colliding) to “compliance preservation” (not violating regulatory limitations).

\section{State of the Art Review}

In the literature on autonomous UAS, vision -based navigation has emerged as a predominant paradigm for enabling perception-driven localization, obstacle awareness, and motion planning. Although these techniques are often motivated by scenarios where Global Navigation Satellite Systems (GNSS) signals are unavailable, unreliable or degraded, numerous studies has demonstrated that vision-based approaches are not limited to such conditions, but rather constitute robust and versatile solution for UAS navigation and trajectory planning across a wide range of operational environments.

 Arafat et al. \cite{arafat2023vision} provide a comprehensive taxonomy of these methods, categorizing vision-based UAS navigation systems, classifying them into three main categories: mapping-based methods for visual localization, object detection and obstacle avoidance approaches, and path planning-based techniques, which directly integrates visual perception into the trajectory generation process. This categorization is aligned with other reviews that also describe vision-based navigation has a multi-stage framework where image acquisition, environmental representation, and motion decision processes are tightly coupled yet conceptually distinct components of autonomous flight systems \cite{huang2025investigation}.

 Inside path planning-based techniques, there is a classification that distinguishes between global and local planning methods. 
 
 The global path planning methods need to have a previous knowledge about the environment, so they present a static solution where the path is previously calculated based on already known factors. 

There are some traditional local path planning methods as graph-based search methods (Dijkstra), that can work properly in static environments with complete maps. But there are significative limitations using these methods for applied solutions as power-line inspection, where it is necessary to have a real time when initially unforeseen obstacles or complex kinematic restrictions, due to the fact that they depend on a previous environment representation. The computational complexity of these methods also increases fast as the operational scenario grows, adding an additional barrier in these dynamic and fast response scenarios \cite{debnath2024review} \cite{meng2025advances}. 

 The local path planning algorithm generates a path while the vehicle is moving through the environment, having the capability of producing a new path based on environmental elements that it encounters during maneuvers. The proposed path planning solution for power-line inspection belongs to this category.
 
 To understand the magnitude of the opportunity that the proposed Local Reactive Model represents, it is imperative to tackle first the operational environment and the severe limitations that the actual literature perspectives are not able to mitigate. Power-line inspection is not a vacuum real-robotic problem; it is a legal engineering operation executed in real time. 

A more specific literature review focused on path planning for Unmanned Aerial Systems \C{referencia Xabi [1]} reveals that there is a systematic limitation in the academic solutions: they are 'toy problem' dependent. In these experimental scenarios, certain assumptions are introduced to reduce the mathematical complexity; however, they undermine the method's applicability in industrial settings: 

\begin{itemize}
    \item Obstacle homogeneity: Obstacles are typically modeled as primitive geometry (spheres, cylinders, cubes) or generic point clouds.There is not clear semantic distinction between a three (where the UAS can approach) or a human (that owns legal aerial exclusion  rights) \C{referencia Xabi [3]}.
    \item Binary Safety: Conventional algorithms define exit as $Distance > 0$. Under this assumption, a mission is deemed successful as long as the drone does not physically impact an obstacle. In practice, however, civil UAS operations are conducted within fully regulated environments, structured into predefined categories and operational scenarios that impose explicit operational limits and constraints. As a result, this assumption becomes inapplicable: Within such regulatory frameworks, a 'near collision' \cite{international2016annex} or the overflight of prohibited elements (such as humans) constitutes a regulatory violation, with legal consequences in terms of administrative enforcement and potential civil liability, regardless of the absence of physical contact \cite{EU2019_947}.  
    
    \item Free Space Assumption: It is assumed that every free air volume not occupied by physical mass is valid to be flown \C{referencia Xabi [2]}. As depicted in previous item, this assumption is in practice not regulatory compliant, as there are prohibited regions defined by law, defined as dynamic exclusion zones. 
\end{itemize}

%Idea general a conseguir: Después de haber introducido los dos tipos de métodos de planificación (global y local) comentar el tema de las limitaciones de los enfoques actuales introduciendo la falacia del toy problem. Luego meterse en el análisis del estado del arte hecho por Xabi y con alguna aportación adicional de las que tengo yo que pueda ser incluida. 
% Después, reformular la sección del SORA de Xabi, para incluir el concepto de persona no involucrada y las diferentes implicaciones que tiene su "esquive" tanto en categoría abierta como en categoría específica con el marco SORA 2.5 y la inclusión de la capacidad de mitigación en tierra con el esquive de personas. 

% Por último, ver cómo tratar la comparativa estratégica, para incluir más métodos y estructurar la tabla de manera un poco más enriquecida. 

\subsection{Regulatory Context: EASA Rules and SORA 2.5 Methodology}

To better understand the potenciality of the proposed solution, it is necessary to explain the regulatory context around the industrial inspection environment. Depending on operational risk, it can be endorsed in different existing categories. Inside all these categories, there are key concepts that transform navigation in an optimization problem with hard restrictions. 

\subsubsection{Open Category and 'uninvolved person' concept }

The open category for UAS constitutes the European regulatory framework established by EASA for low-risk UAS operations. In this category, no prior operational authorization or operational declaration is required, provided a set of predefined limitations that has to be respected \cite{EU2019_947}. Among these limitations, there is one concept directly related to the proposed solution: the proximity to 'uninvolved person'. 

The concept of 'uninvolved person is defined by EASA as a person who is not participating in the UAS operation or who is not aware of the instructions and safety precautions given by the UAS operator. In contrast to a power-line tower or a cable, an 'uninvolved person' generates a dynamic risk field in the surrounding space. 

Within the open category, excluding the A1 subcategory (not scalable for industrial inspections), the permitted proximity to 'uninvolved person' constitutes a key criterion used to define the limitations and the associated risk mitigation. As a generalization, the open category is structured according to a graduated risk-based approach whereby decreasing separation distances from 'uninvolved person' are compensated by progressively stricter aircraft design constraints and higher pilot competency or strict path definitions. 

Within this risk-based approach, there are two rules that are going to be used as a reference for the definition of the proposed solution: 

\begin{itemize}

    \item 1:1 Rule: This rule establishes that, the UAS has to maintain a horizontal separation to 'uninvolved people' that has to be greater or equal to the flight height.

    \item dynamic implications: In general, and if there is no other specific limitation or constraint, if the UAS ascends, for example, from 20 to 50 meters to inspect upper part of a power-line asset, the surrounding exclusion radius from 'uninvolved person' on ground increases instantaneously from 20 to 50 meters. A static path planning model is not capable to manage this variable geometry. 
    
\end{itemize}


\subsubsection{Specific Category and Ground Risk Buffer Concept }
%adaptación de la descripción de Xabi para darle un contexto un poco más regulatorio y menos definición genérica. 
The specific category fo UAS constitutes the European regulatory framework established by EASA for operations whose risk level exceeds the predefined limitations of the open category, but can be demonstrated as acceptable through a Specific Operation Risk Assessment (SORA) and the applied appropriated mitigations \cite{EU2019_947}.

Within this category, the protection of uninvolved persons remains a central regulatory objective; however, it is no longer implemented through fixed operational separation distances, as in the open category. Instead, risk to 'uninvolved people' on the ground is addressed through the definition and control of a ground risk buffer, understood as the spatial margin surrounding the intended Operational Volume. This buffer explicitly accounts for the potential impact area resulting from a catastrophic failure of the UAS, including loss of control or uncontrolled descent, and is dimensioned to ensure that the probability and severity of harm to 'uninvolved people' remain within acceptable limits. The Ground Risk Buffer is therefore managed through a combination of operational constraints, technical mitigations and procedural measures, forming an integral element of the overall safety management. Within this Ground Risk Buffer definition, the Kinetic Energy ($v, m$) and height ($h$) of the UAS represent important limitation parameters. 

The actual path-planning algorithms optimize the velocity to reduce the mission time ($J{tiempo}$). However, as the velocity increases, the UAS expands the Ground Risk Buffer, according to the formulas of SORA Annex F \cite{JARUS_SORA_2_5_2024}. A standard planner from the literature will accelerate to improve efficiency \cite{dong2017efficient}, inadvertently violating the legally required safety margin.


\subsubsection{Strategic Mitigation to reduce the Ground Risk Class }

The new risk analysis methodology, SORA 2.5, presented by JARUS \cite{JARUS_SORA_2_5_2024} and adopted by EASA establishes a set of ground risk mitigation factors that can be applied to reduce the ground risk level of the operation before flight, for Unmanned Aerial Systems. This reduction in ground risk level is crucial, as it will imply a reduction in the intrinsic risk of a person being stuck by an UA and a reduction in the level of risk of the whole operation and will allow the aircraft to fly over the expected area without the need to pass a certification process. 

Inside these strategic mitigations, the 'M1 (C) mitigation-Ground observation' is directly related with the concept of 'uninvolved people'. This mitigation implies that the UAS crew or the systems inside the aircraft itself can observe most of the overflown area, allowing the detection of 'uninvolved people' in the operational area and maneuvering the UA so that the 'uninvolved people' are not overflown \cite{EASA_AMC_GM_SORA25_2025}. The proposed solution addresses this mitigation, proposing an innovative solution with a real time path-planning tool that modifies the planned route to evade these 'uninvolved people'. 

%Falta la parte de estudio comparativo del estado del arte hecho por Xabi e integrando algún ejemplo más de los encontrados por mi. Además incluir al final la tabla comparativa entre nuestro modelo, los modelos matemáticos y las soluciones comerciales. 


%\section{Methodology}
\section{Proposed Solution: Local Reactive Model for Real Time UAS Navigation}

This section is focused on the flow description of the proposed solution. As the model is designed for an UAS platform, it has to take into account every step involved in the flight, from the mission reading and takeoff of the aircraft, passing through the guided flight and conflict resolution (static and dynamic obstacle approach), and finally the landing phase, once the mission is finished or caused by any problem encountered throughout the flight. 

The objective is to complete a powerline-inspection in an automatic way following a Flight Plan. But the special characteristic of this Flight Plan is that it is dynamic, and if any regulatory barrier or ordinary obstacle item (as a powerline asset) appears in the expected trajectory of the UAS, it is capable of performing a Flight Plan reconfiguration. 

The initial mission Flight Plan was generated based on the information available during the pre-flight phase, using the global planning tools and methodologies. In particular, the planning process relied on the nominal positions of the power-line assets, which may be outdated, incomplete, or inaccurately geo-rreferenced. The Flight plan was defined using a map-based planning approach that incorporates additional constraints and contextual information, including communication network coverage, the availability of external services (e.g., GNSS or LTE), and previously known obstacle locations. These factors collectively influence the definition of Flight Plan waypoints (WPs),

Up to this point, all strategic mitigation measures that can be evaluated and applied prior to the operation during the pre-flight phase have been systematically addressed in alignment with the corresponding operational risk assessment. These measures ensure compliance with the applicable European regulatory framework and consistency with the target operational category. 

However, as introduced in Section 2, the operational focus shifts during the execution phase. At this stage, residual risks may arise due to dynamic, uncertain, or locally perceived conditions that cannot be fully captured during the planning phase. In these circumstances, simplified assumptions and toy-problem formulations often emerge, together with a neglect of regulatory compliance considerations.The proposed local reactive navigation model enables realistic real-time adaptation to the environment. The functional workflow of the proposed solution is illustrated in the following figure: 


% Dentro de los estados de vuelo tenemos: "Navegando, Obstáculo detectado, Evadiendo con A*)


\begin{lstlisting}[language=Python]

WAYPOINTS_FILE, EARTH RADIUS_M, SAFETY_DISTANCE_M, DETECTION_RANGE_M 
REACTION_DISTANCE_M,ARRIVAL_TOLERANCE_M, ALTITUDE_TOLERANCE_M, 
EVASION_VELOCITY_LATERAL_MS, EVASION_SPEED_MS.
WP_TOLERANCE_M = ARRIVAL_TOLERANCE_M

\end{lstlisting}

% estas serán los valores que definirán el estado en cada instante de la misión. En función de los parámetros de la variable estado se incurrirá en unas acciones u otras. 

FlightPlan <- load_mission()
FlightMode = ['GUIDED']
state ['Telemetry'] <- get_telemetry() %obtener la telemetría de la aeronave.
tel = state['Telemetry'] %asignar parámetro tel al estado de la telemetría.
state['obstacles'] <- get_detection() %activar la detección para buscar los obstáculos durante el vuelo, Falta definir lo que va dentro de esta función y cómo relacionarla con el bucle de control. => generar bucle de detección corriendo todo el rato.     
%elemento clave, primero comprueba que no hay obstáculos en cada bucle y luego si no hay obstáculos pasa por el código de navegación normal, de manera que mientras que haya obstáculo, navega por el path de conflicto y no por el calculado con anterioridad. Esto resulta clave, puesto que soluciona el problema y después vuelve lo más rápido posible a su path previo, de manera que se asegura la resolución del conflicto pero la vuelta lo más rápido posible al path original, pensado para optimizar la captura de datos, o el objetivo que se tuviera en concreto (buscado en la planificación global inicial). 
        
while FlightMode != ['LAND']
    %cada vez que se reinicia el bucle del while se regeneran en función del estado los parámetros necesarios para las acciones
    state ['Telemetry'] <- get_telemetry()
    state['obstacles'] <- get_detection()
    tel = state[`telemetry`]
    obs = list(state['obstacles'])
    current_idx = state['current_wp_idx']
    wps = state['waypoints'] %¿es lo mismo que mi variable Flight Plan?
    active_path = [] %siempre limpio el parámetro para asegurarme que no se queda grabado nada que no quiera
    path_idx = 0 %lo mismo que el parámetro anterior, para volver a dejar el parámetro limpio antes de reiniciar el bucle

    %chequeo para comprobar si activar el camino de evasión 
    active_path = state['evasion_path'] %cuando se haya activado el evasion path se activará para esquivar el problema. 
    path_idx = state['path_index']
    if not active_path: 
        nearest_obs = None
        min_dist = float('inf')
        for o in obs:
            d = o.get('distance', 9999)
            if d < min_dist:
                min_dist = d
                nearest_obs = o
        %aquí se produce el cambio clave cuando ha aparecido un objeto y se alcanza la distancia mínima de reacción
        if nearest_obs and min_dist < REACTION_DISTANCE_M:
            %obstáculo dectectado a distancia min_dist. Se comienza a Planificar ruta A*.
            target_wp= wps[current_idx] if current_idx < len(wps) else wps[-1]
            new_route = planner.plan_route(tel['lat'], tel['lon'], target_wp['lat'], target_wp['lon'], obs)
            if new_route:
                state['evasion_path'] = new_route
                state['evasion_grid_origin'] = {'lat': tel['lat']}, 'lon': tel['lon']} 
                state['path_idx']= 0
                state['evasion_active'] = True
                active_path = new_route
            else 
                %No ha funcionado la búsqueda de una ruta alternativa con A*. Hay riesgo de colisión. 

    if active_path: %Condición que será verdadera mientras que haya una evasión en marcha. 
        if path_idx < len(active_path) %mientras que se esté calculando un camino alternativo y no se llegue al último punto sigo el en bucle de evasión.
            ...
            continue %para volver al inicio del while sin continuar hasta abajo
        else
            state['evasion_path'] = [] %como ya se ha alcanzado el último punto del active path, el camino de evasión se vacía.
            state['evasion_active] = False %se desactiva la evasión para volver a la navegación normal hasta que haya otro obstaculo.

    %navegación estándar 
    if current_idx < len(FlightPlan)
        target = FlightPlan [current_idx] 
        dist = haversine(tel['lat], tel['lon'], target['lat'], target['lon'])
        if dist < WP_TOLERANCE_M:
            state

    else %si se ha alcanzado el último wp del FlightPlan, se pone modo aterrizaje y se sale del bucle de control
        tel['armed'] = False
        FlightMode = ['LAND']



def load_mission ():
def get_telemetry(): 
def rx_obstacles(): 
%aplicación corriendo en paralelo para la detección de obstáculos. Si detecta un obstáculo le meterá al sistema su posición y el algoritmo de evasión se encargará de calcular si la distancia está dentro de los márgenes críticos. De ser así entrará el juego el porce_manager.py (ver explicación más abajo). 
def status(): 
%para saber el estado en un json en el que se encuentra la misión 
def mavlink_loop(): 
%simulador corriendo en paralelo con el código de navegación. Le llegarán al simulador los cambios de ruta y de estado que vaya generando el control loop. 
def control_loop():
% código principal de la solución. En él mientras que se haya cargado una misión, comienza con su lectura y analiza el estado del vuelo. Una vez que ha despegado y se encuentra en ruta comienza con el algoritmo de navegación: 
%- Mientras que no se detecte ningún obstáculo, la navegación es estándar ( línea 373 a 394 del flight_controller.py.  Si se alcanza el último waypoint que había programado, el sistema comanda el modo LAND y termina la misión. 
%- En el caso de durante el transcurso de la misión aparezca un obstáculo, se va a la parte del código de algoritmo porce (evasión) lineas 316 a 373. Al aparecer un obstáculo la variable nearest_obs es true y lo que se mira es si min_dist (que es la distancia al obstáculo más cercano) está dentro del umbral de distancia de reacción. De ser así se arranca el planificador de ruta alternativa (porce_manager.py). Aquí aparece el grid y se empieza a esquivar el obstáculo y a seguir la ruta alternativa que se va calculando mientras que el obstáculo esté presente y esté dentro de la distancia de reacción. Mientras que exista riesgo de colisión se navega por el path de evasión, hasta que la evasión se completa y lo que ocurre es que se vuelve al estado de vuelo normal con el path que quede por recorrer (linea 373 a 394). 

%Va a ser necesario reformular teóricamente el código de Xabi en este pseudocódigo, para que tenga sentido sin correr una app por detrás y sin estar bebiendo de mavlink. Dar u

In the first step of the proposed method the $load_plan()$ function read the predefined route WPs, and transfers it to the command unit of the UAS, providing it what is needed to develop the flight in an automatic way. 

Next step is to start reading the telemetry state of the UAS with the $get_telemetry()$ function, to assign to the telemetry state variables inside the $state$ array its initial values. 
The $tel$ variable will be assigned to $state['telemetry']$ for simplicity along the proposed solution. 

The state variables are stored inside an array with initial values as follows: 
\begin{lstlisting}[language=Python]
state = {
    'telemetry': {
        'lat': 0.0, 'lon': 0.0, 'alt': 0.0,
        'roll': 0, 'pitch': 0, 'yaw': 0,
        'heading': 0,
        'armed': False,
        'mode': 'UNKNOWN',
        'last_update': 0
    },
    'waypoints': [],
    'current_wp_idx': 1,
    'obstacles': [],
    'last_obstacle_update': 0,
    'evasion_active': False,
    'evasion_path': [],
    'evasion_grid_origin': None, 
    'path_index': 0,}
\end{lstlisting}
%Pendiente definir el estilo para que los códigos de Python queden con color y fondo como una figura
It is assumed that the UAS is configured by default to start in ['GUIDED'] mode, and that the aircraft it is already armed. 

At this point is where the control loop starts. While the flight plan is not completed and, therefore, entering on ['LAND'] mode ($while FlightMode != ['LAND']$, the proposed solution will run continuously, checking to see whether it is necessary or not to redefine the planned path based on events that occur during the mission. 
Each time that this control loop starts, the state variables ['Telemetry'], ['obstacles'] ['current_wp_idx'], ['waypoints']... are regenerated depending on what happen in the previous execution of the loop. %Esto se ha cambiado por la función referesh_telemetry_and_obstacles()

%Meter explicación de que implica el if de telemetry stale

%en este punto se ejecuta el chequeo de si hay alguna acción de failsafe. Luego más adelante en el pseudocódigo se ve cuándo se activa el failsafe. Esto se sitúa aquí simplemente porque en el caso de que haya una acción de failsafe (como en el caso crítico el land) se debe ejecutar lo primero de todo antes de seguir analizando esa iteración del loop.



Inside this control loop, the detection function is activated, so the system will continuously check the obstacles state with the $get_detection()$ function. The system responsible for providing the input to this function is an RGB camera on-boarded in the UAS. The methodology used for detection of target non regulatory compliant items is based on the tracking proposed in \cite{alaez2023real}.
%la función get_detection se llama ahora nearest y hace básicamente lo mismo, chequear el estado de los obstáculos (devolviendo 0 en el caso de que no haya ninguno en el rango de detección de la cámara). Si hay algún obstáculo se queda con el más cercano. The system responsible for providing the input to this function is an RGB camera on-boarded in the UAS. The methodology used for detection of target non regulatory compliant items is based on the tracking proposed in \cite{alaez2023real}.

This detection function is the key element to select the path inside the control loop. If no conflict is detected, the system executes the nominal navigation flow, tracking the previously planned trajectory. This behavior is represented in the figure, between lines %insertar líneas

The nominal navigation works as follows: 
%En el nuevo pseudocódigo esta navegación nominal entra en el código a partir de la línea 24, cuand o ha chequeado que la evasión no está activa y entonces se siguen los puntos de la misión. Se mete de nuevas la parte de chequear la distancia de tolerancia con el current_wp. Si está demasiado cerca se le manda al siguiente. 

- If $current_idx$ is less than the length of the proposed flight plan, the system continues following the planned flight path. At each iteration, it checks whether the distance to the current target waypoint is below $WP\_TOLERANCE$; if so, the target is updated to the next waypoint for the subsequent iteration.
%Se mete de nuevas en el pseudocódigo de la web de Xabi esta parte de chequear la distancia de tolerancia con el current_wp. Si está demasiado cerca se le manda al siguiente. 

- If not, it implies that the last waypoint of the Flight Plan is reached, so the Flight Mode will be switched to $['LAND']$. This action will cause the close of the control loop.
%metido en el else del pseudocódigo. 

%propuesta de darle la vuelta y primero explicar la evasión cuando la navegación ha sido inhibida.
Conversely, when an obstacle is identified, nominal navigation is temporarily inhibited and the system switches to a conflict-resolution mode, in which the vehicle progresses along an alternative avoidance path.

%meter aquí la explicación de dentro del if de nearest. Se comprueba si hay obstáculo (nearest =True), Se compara la distancia de reacción dinámica con la distancia a a que se encuentra el obstáculo a ver si la distancia a la que está el obstáculo es menor a la de reacción, y se compara si es posible replanificar en ese momento (se necesita un buffer de tiempo entre replanificación de rutas para no saturar a la aeronave). 

%Aquí falta explicar que es en este punto donde a través de [route] se calcula la ruta de evasión a través de la solución porce (llamada astar_local en el pseudocódigo). Comentar que se describe más adelante esta solución/función en detalle. Si la ruta es válida, es decir tiene más de dos puntos se activa la evasión. Si no se aplica el failsafe que toque. 

This behavior is essential because it guarantees real-time conflict resolution and, once safe conditions are restored, it drives the system back to the original plan as quickly as possible. In this way, operational safety is maintained without compromising the efficiency objective defined during the initial global planning stage, whether that objective is optimized data acquisition or any other mission-specific metric established a priori.











%\section{Algorithm Design: }
\section{Path planning and Obstacle avoidance Real-time Collision Evasion (PORCE) model}

In this section, the PORCE (Path planning and Obstacle avoidance Real-Time Collision Evasion) model is formally defined. The navigation problem is formulated as a trajectory optimization task subject to strict safety constraints in a partially observable environment. The theoretical justification of a hybrid architecture (Deep Perception + Deterministic Planning) is presented, in contrast to purely stochastic methods, for safety-critical aerospace systems.

\subsection{Mathematical Formulation}
\subsubsection{State Space and Environment}

Let $\mathcal{W} \subset \mathbb{R}^3$ denote the global operational space. Under the operational assumption of an \textit{Infinite Column} (designed to avoid cables and vertical structures), the problem is projected onto the $\mathcal{W}_{2D} \subset \mathbb{R}^2$ plane.

The vehicle state at the discrete time $t$ is defined as vector: 

\begin{equation}
    \mathbf{x}_t = [p_t, v_t, \psi_t]^T \in \mathbb{R}^2 \times \mathbb{R}^2 \times [-\pi, \pi),
\end{equation}

where: 
\begin{itemize}
    \item $p_t = (x, y)$ denotes the cartesian position within the local frame. 
    \item $v_t$ is the inertial velocity vector.
    \item $\psi_t$ is the orientation angle (\textit{yaw}), coupled to the velocity vector such that $\psi_t = \arctan(v_y / v_x)$.
\end{itemize}

\subsubsection{Obstacle Modeling and Safety Regions}

The perception system (YoloV8 based) detects a set of obstacles $\mathcal{O}_t = \{o_1, \dots, o_k\}$ during time $t$. Each obstacle is treated as a stochastic entity that has to be bounded. 
An Augmented Safety Region is defined $\mathcal{Z}_{safe}(o_i)$ for an obstacle $i$ centered at $p_{obs}$ with a safety radius $R_s \text{m}$ as: %quitado el valor numérico del safety radius. Luego darlo en la sección de resultados si se quiere.

\begin{equation}
    \mathcal{Z}_{safe}(o_i) = \{ q \in \mathcal{W}_{2D} : \| q - p_{obs} \|_2 \le R_s \}
\end{equation}

As a consequence, flyable free space $\mathcal{F}$ where the UAS is allowed to operate is defined as the subtraction of all prohibited regions from the total space:

\begin{equation}
    \mathcal{F} = \mathcal{W}_{2D} \setminus \bigcup_{i=1}^{k} \mathcal{Z}_{safe}(o_i)
\end{equation}

\subsubsection{Discrete Planning (Grid Search Optimization)}

To solve the real time path problem, the local environment is discretized in an occupancy grid $\mathcal{G}$ of dimension $L \times L$ with spatial resolution $\delta$. %en vez de 400 he puesto NxN para que sea genérico

Each cell  $c_{ij} \in \mathcal{G}$ of the occupancy grid has a binary state $S(c_{ij})$ defined as: 

\begin{equation}
    S(c_{ij}) = 
    \begin{cases} 
    1 \quad (\text{Occupied}) & \text{si } \exists p \in c_{ij} \text{such as } p \notin \mathcal{F} \\ 
    0 \quad (\text{Free}) & \text{other cases} 
    \end{cases}
\end{equation}

The cost function used by the $A*$ Algorithm to evaluate the candidate node $n$ is the following: 

\begin{equation}
    f(n) = g(n) + \epsilon \cdot h(n),
\end{equation}

where:
\begin{itemize}
    \item $g(n)$ is the accumulated real cost from the initial node $n_{start}$ to $n$.
    \item $h(n) = \| p_n - p_{goal} \|_2$ is the admissible Euclidean heuristic towards the objective. 
    \item $\epsilon \ge 1$ is a ponderization factor (configured to guarantee strict optimality).
\end{itemize}

\subsubsection{Triggering Logic}
The controller continuously supervises an imminent risk function $\mathcal{R}(t)$. The transition between \textit{Mission} mode and \textit{Evasion} mode is activated if and only if:

\begin{equation}
    \exists o_i \in \mathcal{O}_t \quad \text{such as} \quad (\| p_t - p_{o_i} \|_2 \le D_{react}) \land (p_{o_i} \in \mathcal{V}_{cone}(v_t)),
\end{equation}

where $\mathcal{V}_{cone}$ represents the projected UAS velocity cone and $D_{react}$ is the reaction threshold. %de nuevo el valor numérico se ha eliminado en esta descripción. Luego para los resultados se pueden comentar los valores presentados.

\subsection{Scientific justification of the approach}

\subsubsection{Determinism vs. Stochastic black boxes}

To introduce autonomous Unmanned Aerial Systems into the market, the regulation states that it  is mandatory to have transparency and traceability of all decisions \cite{EASA_AI_Roadmap_2.0}. In the End to End (DRL) solutions, a $\pi_\theta(a|s)$ politics totally learned by reinforced learning, is a probabilistic "black box". It is not possible to guaranty that it will not fail in the edge cases. 

However, in the PORCE model, the A* algorithm provides a mathematical guaranty of completeness. If a geometrically valid path exists $\sigma \subset \mathcal{F}$, the algorithm will always find it . 

\subsubsection{Edge Computing Feasibility}

The temporal complexity of the A* algorithm in the worst case scenario $O(b^d)$, where $b$ is the ramification factor and $d$ the deepness. As the searching horizon is limited to a bounded local grid ($LxL$) and a defined resolution $\delta$, the state space is finite and small (e.g., for $L =400$ m and $r=10$ m the state space is $40 \times 40$ cells).

This limited state space guaranties that the computational time $T_{comp}$ meets: 

\begin{equation}
    T_{comp} \ll T_{loop} \approx 100\text{ms},
\end{equation} 

which validates the feasibility of the system to operate on resource-constrained onboard hardware (e.g., Jetson or Raspberry Pi platforms) while maintaining a safe control frequency at typical cruise speeds of $8\text{m/s}$. 

\subsubsection{Geometric Constraints Handling}

The enforcement of the {Infinite Column Assumption} is a hard constraint injected into the configuration space. In contrast to Reinforced Learning, which must learn to avoid collisions through negative penalties (rewards), the PORCE approach makes it geometrically impossible to plan a route through an obstacle, thereby guaranteeing \textit{Zero-Shot} safety, i.e. without the necessity of previous training in the same specific scenario. 

%Se ha hecho una explicación matemática del algoritmo Porce, pero puede ser interesante una explicación de un pseudocódigo más acortado como se ha hecho en la sección anterior con el hilo principal de código. 

\section{Implementation and Results}

%Se empieza mostrando una gráfica del camino definido para el dron. Con el instante inicial del gif de Xabi valdría. En ese instante se ven el UAS y la ruta de WP que tiene que seguir. Para probar la eficacia del código se hace que la posición de las torres no sea conocida por el UAS, de manera que durante el transcurso de la misión también tendrá que esquivarlas. 

%La siguiente gráfica que se muestra es el primer instante en el que detecta un obstáculo pero todavía no tiene que calcular el path alternativo. De esta manera, el obstáculo se representa en la imagen, puesto que el algoritmo de detección lo ha captado y ha calculado su distancia a este. 

%La siguiente gráfica muestra el instante en el que el obstáculo está dentr ode la reaction distance y el sistema comienza a calcular la ruta (se pinta en naranja la ruta propuesta por el UAS). 

%La siguiente imagen muestra la resolución del conflicto completa, una vez el UAS ha alcanzado el camino inicial de nuevo, para poder seguir con el objetivo de la captura de datos. En ese punto quedará reportado el problema (y en el análisis futuro puede que falten datos del activo, pero la seguridad prima sobre la obtención de los datos). 

% Por último se muestra la imagen final del gif completo, dónde se ve como el sistema ha sido capáz de ir resolviendo los diferentes conflictos que se ha ido encontrando a lo largo de la misión.


% POSIBLES COSAS A METER: "Problema local de la evasión de un obstáculo en movimiento, en el que se muestre como se reacciona ante esto y cómo evoluciona el cálculo del camino de solución propuesto. 


\section{Conclusion}

\section{Future works}


\printbibliography

\end{document}